% Decoder block information flow and one causal-attention read.
% Equal-sized cells encode attention weights by shade and number, not width.
\documentclass[border=2mm]{standalone}
\usepackage{amsmath,amssymb}
% Chinese labels in figures; fonts are shared with the book.
\usepackage[UTF8,fontset=none]{ctex}
% CJK fonts for the book, the figures and the slides.
% Noto Sans CJK SC (SIL Open Font License) is the same typeface as Source Han Sans SC.
% macOS: brew install --cask font-noto-sans-cjk-sc
% Debian/Ubuntu: apt install fonts-noto-cjk
\setCJKmainfont[BoldFont={Noto Sans CJK SC Bold}]{Noto Sans CJK SC}
\setCJKsansfont[BoldFont={Noto Sans CJK SC Bold}]{Noto Sans CJK SC}
\setCJKmonofont{Noto Sans CJK SC}


% Shared visual language for the figures of the book.
% Every figure in this directory loads this file so the chapter reads as one system.
% NOTE: no blank lines are allowed inside \tikzset{...} (a blank line emits \par
% and breaks key scanning).
\NeedsTeXFormat{LaTeX2e}
\ProvidesPackage{figstyle}[2026/10/21 Figure style for Free Energy to LLMs]

\RequirePackage{tikz}
\usetikzlibrary{arrows.meta,positioning,calc,fit,shapes.geometric,shapes.misc,%
  backgrounds,matrix,decorations.pathreplacing,decorations.markings,patterns.meta}

% ---------------------------------------------------------------- palette
% Same values as book/preamble/style.tex, prefixed with F.
\definecolor{FNavy}{HTML}{14324A}
\definecolor{FTeal}{HTML}{0E7490}
\definecolor{FSky}{HTML}{E7F5F8}
\definecolor{FOrange}{HTML}{D97706}
\definecolor{FAmber}{HTML}{FFF3D6}
\definecolor{FInk}{HTML}{1F2937}
\definecolor{FSlate}{HTML}{64748B}
\definecolor{FMist}{HTML}{F5F7FA}
\definecolor{FRose}{HTML}{9F1239}
\definecolor{FRoseMist}{HTML}{FFF1F2}

% ---------------------------------------------------------------- semantics
% Teal   = probability / model / data object
% Orange = training objective, cost, or the quantity moved through the system
% Navy   = neutral structure
% Slate  = secondary label or metadata
% Rose   = failure mode, hazard, masked-out region
\tikzset{
  % Lay figures out in millimetres.  A figure designed 168mm wide is included at
  % \textwidth, so its fonts appear at exactly the size written here.
  every picture/.style={x=1mm,y=1mm,font=\small,line cap=round,line join=round},
  % ---- block styles
  nbase/.style={draw,rounded corners=1.2pt,align=center,inner sep=2.4pt,%
    minimum height=7mm,text=FInk},
  nteal/.style={nbase,draw=FTeal,fill=FSky},
  norange/.style={nbase,draw=FOrange,fill=FAmber},
  nnavy/.style={nbase,draw=FNavy,fill=white},
  nsoft/.style={nbase,draw=FSlate!55,fill=FMist},
  nrose/.style={nbase,draw=FRose,fill=FRoseMist},
  ghost/.style={draw=none,fill=none,align=center,inner sep=1pt},
  plain/.style={draw=none,fill=none,align=center,inner sep=0pt},
  % ---- arrows
  ar/.style={-{Latex[length=1.7mm,width=1.35mm]},draw=FSlate,line width=0.75pt},
  ard/.style={ar,dashed},
  arteal/.style={ar,draw=FTeal},
  arorange/.style={ar,draw=FOrange},
  arrose/.style={ar,draw=FRose},
  thickar/.style={ar,line width=1.2pt},
  % ---- text roles
  ftitle/.style={font=\large\bfseries,text=FNavy,align=center},
  ptitle/.style={font=\small\bfseries,text=FNavy,align=center},
  ptitleo/.style={font=\small\bfseries,text=FOrange,align=center},
  lbl/.style={font=\scriptsize,text=FSlate,align=center,inner sep=1pt},
  lblt/.style={lbl,text=FTeal},
  lblo/.style={lbl,text=FOrange},
  lbln/.style={lbl,text=FNavy},
  lblr/.style={lbl,text=FRose},
  sub/.style={font=\scriptsize,text=FSlate!90,align=center,inner sep=1pt},
  monob/.style={font=\ttfamily\scriptsize,text=FInk,inner sep=1pt},
  mathb/.style={font=\small,text=FInk,inner sep=1pt},
  eqn/.style={font=\scriptsize,inner sep=1.5pt,text=FInk},
  % ---- decorations
  brace/.style={decorate,decoration={brace,amplitude=2.4pt},draw=FSlate,%
    line width=0.6pt},
  bracet/.style={brace,draw=FTeal},
  braceo/.style={brace,draw=FOrange},
  axis/.style={draw=FSlate!55,line width=0.6pt},
  guide/.style={draw=FSlate!45,line width=0.5pt,dashed},
  halolabel/.style={font=\scriptsize,text=FInk,fill=white,fill opacity=0.85,%
    text opacity=1,inner sep=1.4pt,align=center},
}

% Figure title above the drawing, anchored at the figure's top-left.
\newcommand{\figtitle}[2][0]{%
  \node[ftitle,anchor=north west,outer sep=0pt] at (#1,0) {#2};}

% Rounded panel behind a subgroup of nodes.
\newcommand{\figpanel}[3][]{%
  \node[draw=FSlate!35,fill=#3,rounded corners=2pt,inner sep=2.6mm,#1] {#2};}

\endinput

\begin{document}
\begin{tikzpicture}
  % One Pre-Norm decoder block. The vertical line is the residual path.
  \node[ptitle,anchor=west] at (0,83) {Decoder block：沿残差流更新};
  \node[nnavy,minimum width=22mm,minimum height=8mm] (hin) at (15,6) {$h^{(l)}$};
  \node[nteal,minimum width=24mm,minimum height=7mm] (norm1) at (43,14) {RMSNorm};
  \node[norange,minimum width=35mm,minimum height=9mm] (attn) at (43,25) {因果多头注意力};
  \node[circle,draw=FTeal,fill=white,line width=.8pt,inner sep=0pt,
    minimum size=5mm,font=\small,text=FTeal] (add1) at (15,36) {$+$};
  \node[nteal,minimum width=24mm,minimum height=7mm] (norm2) at (43,43) {RMSNorm};
  \node[nteal,minimum width=35mm,minimum height=9mm] (ffn) at (43,55) {SwiGLU / FFN};
  \node[circle,draw=FTeal,fill=white,line width=.8pt,inner sep=0pt,
    minimum size=5mm,font=\small,text=FTeal] (add2) at (15,66) {$+$};
  \node[nnavy,minimum width=22mm,minimum height=8mm] (hout) at (15,76) {$h^{(l+1)}$};

  \draw[arteal,line width=1.1pt] (15,10)--(add1.south);
  \draw[arteal,line width=1.1pt] (add1.north)--(add2.south);
  \draw[arteal,line width=1.1pt] (add2.north)--(hout.south);
  \node[lblt,rotate=90] at (3,42) {残差路径};
  \draw[ar] (15,14)--(norm1.west);
  \draw[ar] (norm1.north)--(attn.south);
  \draw[ar] (attn.north)--(43,36)--(add1.east);
  \draw[ar] (15,43)--(norm2.west);
  \draw[ar] (norm2.north)--(ffn.south);
  \draw[ar] (ffn.north)--(43,66)--(add2.east);

  % Causal attention: rows are queries, columns are keys.
  \node[ptitle,anchor=west] at (76,83) {注意力：位置 3 怎样读取历史};
  \node[mathb] at (122,74)
    {$A=\operatorname{softmax}(QK^{\mathsf T}/\sqrt{d_k}+M),\quad O=AV$};
  \node[lbln,anchor=west] at (128,68) {列：key $j$};
  \foreach \c in {1,...,5} {
    \pgfmathsetmacro{\xx}{96.5+5*\c}
    \node[lbl] at (\xx,67.8) {\c};
  }
  \foreach \r in {1,...,5} {
    \foreach \c in {1,...,5} {
      \pgfmathsetmacro{\xx}{99+5*(\c-1)}
      \pgfmathsetmacro{\yy}{66-5*\r}
      \ifnum\c>\r
        \fill[FRoseMist] (\xx,\yy) rectangle ++(5,5);
      \else
        \fill[FSky] (\xx,\yy) rectangle ++(5,5);
      \fi
    }
  }
  \fill[FTeal!28] (99,51) rectangle (104,56);
  \fill[FTeal!45] (104,51) rectangle (109,56);
  \fill[FTeal!76] (109,51) rectangle (114,56);
  \draw[FSlate!60,line width=.45pt] (99,41) grid[step=5] (124,66);
  \draw[FOrange,line width=1pt] (99,51) rectangle (124,56);
  \node[lbln,anchor=east] at (96,53.5) {query $t=3$};
  \draw[arorange] (96.5,53.5)--(98.5,53.5);
  \node[lblr,anchor=west,align=left] at (128,59)
    {未来位置 $j>t$\\被遮蔽};
  \draw[arorange] (111.5,51)--(111.5,35);
  \node[lblo,anchor=west] at (114,36) {第 3 行 softmax 权重};

  % Every weight is one equally wide cell. The zero cells mask future values.
  \foreach \x/\num/\tone/\j in {84/0.2/28/1,100/0.3/45/2,116/0.5/76/3} {
    \node[draw=FTeal,fill=FTeal!\tone,minimum width=12mm,
      minimum height=8mm,inner sep=0pt,font=\scriptsize,text=FNavy]
      at (\x,27) {\num};
    \node[nteal,minimum width=12mm,minimum height=8mm] at (\x,13) {$v_{\j}$};
    \draw[arteal] (\x,22.7)--(\x,17.4);
  }
  \foreach \x/\j in {132/4,148/5} {
    \node[draw=FRose,fill=FRoseMist,minimum width=12mm,
      minimum height=8mm,inner sep=0pt,font=\scriptsize,text=FRose]
      at (\x,27) {$0$};
    \node[nsoft,minimum width=12mm,minimum height=8mm] at (\x,13) {$v_{\j}$};
    \draw[axis] (\x,22.7)--(\x,17.4);
  }
  \node[mathb] at (116,3) {$o_3=0.2v_1+0.3v_2+0.5v_3$};
\end{tikzpicture}
\end{document}
